\section{Each Text in Its Setting}
\label{sec:each-element}

The liturgical books say how these texts stand to one another, and they say
little. The introduction to the \work{Ordo lectionum Missae} explains that in
Advent the readings are not arranged as on the Sundays of the year, where the
letter and the Gospel are read in course. Each Advent Mass has readings chosen
for the season: on the First Sunday the Gospel concerns the Lord's coming at
the end of time; the Old Testament readings are prophecies of the Messiah and
of the messianic age, chiefly from Isaiah; the readings from an apostle are
exhortations and proclamations suited to the season (\work{Praenotanda}
nos.~67 and~93). The psalm responds to the first reading, and the Alleluia verse
greets the Gospel (\work{General Instruction of the Roman Missal} 61--62). The
\work{Universal Norms on the Liturgical Year} give Advent a double character:
it prepares for the solemnity of Christmas, in which the first coming of the
Son of God is remembered, and it turns minds to his second coming at the end
of time (no.~39). The antiphons, prayers, Preface and blessing are the
Missal's for this Sunday in each of the three years; the readings are the
second year's.

The \work{Ordo} sets a title over each reading, and no.~123 of its
introduction says that a title shows the chief theme of its reading. The three
titles here are \latin{Utinam dirumperes caelos et descenderes!}, ``O that
thou wouldst rend the heavens and come down!'', over the prophet;
\latin{Revelationem exspectamus Domini nostri Iesu Christi}, ``We await the
revelation of our Lord Jesus Christ'', over the letter; and \latin{Vigilate:
nescitis quando dominus domus veniat}, ``Watch: you know not when the lord of
the house comes'', over the Gospel. The second recasts Paul's participle,
spoken of the Corinthians, into the first person plural of those who hear it.
The books pair none of the readings with another.

\subsection{The Gospel: the last words before the Passion}

Mark's thirteenth chapter is the longest speech of Jesus in that Gospel, and
the last before the account of his Passion begins: the next verse is ``Now the
feast of the pasch and of the Azymes was after two days'' (14:1). It opens
when a disciple admires the stones of the Temple and Jesus says that not a
stone will be left on a stone. On the Mount of Olives, opposite the Temple,
Peter, James, John and Andrew ask him privately, ``Tell us, when shall these
things be and what shall be the sign when all these things shall begin to be
fulfilled?'' (13:3--4). He answers with warnings of deceivers, wars and
persecutions, the abomination of desolation, the sun darkened and the Son of
man coming in the clouds with great power and glory, and the fig tree whose
leaves tell that summer is near. Then he refuses the question of the date:
``But of that day or hour no man knoweth, neither the angels in heaven, nor
the Son, but the Father'' (13:32). The reading begins at the next verse.

The reading opens with two commands, \latin{Videte, vigilate} (13:33): look,
keep awake. The verb \latin{vigilare} means to stay awake through the
night, and the night of v.~35 is divided by the four Roman watches: evening,
midnight, cockcrow and morning. The parable is brief and has no ending. A man
going abroad leaves his house, gives his servants authority, assigns to each
his work, and commands the doorkeeper to stay awake; he does not come back
within the story. The last verse breaks out of the private conversation with
four men: ``what I say to you, I say to all'' (13:37).

\subsection{The first reading: a prayer from the ruins}

Isaiah 63:7--64:11 (Douay 64:12) is one long prayer. It begins by recalling the Lord's
mercies to Israel in the Exodus, ``the angel of his presence saved them''
(63:9), and Moses leading them through the deep; it asks where that Lord is
now (63:11, 15); it confesses sin; and it ends before a ruin: ``The city of thy
sanctuary is become a desert \ldots\ The house of our holiness, and of our
glory, where our fathers praised thee, is burnt with fire'' (64:9--10; Douay
64:10--11). The people who speak have lost the sanctuary and say so: ``our
enemies have trodden down thy sanctuary'' (63:18). The reading takes three pieces of the prayer:
the people's appeal to God as Father and Redeemer (63:16b--17), the cry for him
to break the sky and come down (63:19b), and a confession that ends with the
potter and the clay (64:2b--7).

In the Latin, the cry \latin{Utinam dirumperes caelos et descenderes} is
followed by \latin{Descendisti}, ``thou didst come down'' (64:2b). Before the
cry the people ask,
``Why hast thou made us to err?'' (63:17). And in the confession the anger
stands first: ``behold thou art angry, and we have sinned'' (64:4; Douay
64:5).

In the verse of the anger the Latin also says \latin{Occurristi laetanti, et facienti
iustitiam}, ``Thou hast met him that rejoiceth, and doth justice'' (Douay
64:5). The Collect has the same verb, \latin{occurrentes}. In Isaiah God
meets the one who does justice; in the Collect the faithful ask to go out with
just works to meet Christ.

\subsection{The psalm: the Shepherd asked to come}

Psalm 80 (79) is a lament for a defeated people. God is addressed as the
Shepherd of Israel, who leads Joseph like a flock and sits enthroned on the
cherubim, the winged figures over the ark in the Temple, and he is asked to
shine forth before Ephraim, Benjamin and Manasseh.
The psalm then tells a history in a figure: God brought a vine out of Egypt,
cast out the nations and planted it, and it filled the land; now its walls are
broken, passers-by pluck it, and the boar out of the wood lays it waste
(vv.~9--14). Three times a refrain asks God to turn his people and show his
face (vv.~4, 8, 20), each time with a fuller name: ``O God'', ``O God of
hosts'', ``O Lord God of hosts''.

The Lectionary takes three strophes: the plea to the Shepherd to come; the
plea to look down from heaven and visit the vine; and the plea for the one at
God's right hand, with the people's promise not to turn back. The response of the \work{Ordo} asks God to make his
face shine (\latin{illustra}), the word of the \work{Nova Vulgata}, where the
older Vulgate has ``shew'' (\latin{ostende}). And v.~15, in the second strophe,
begins \latin{Deus virtutum, convertere}, ``Turn again, O God of hosts'', the
same imperative that Isaiah's people address to God at 63:17
(\latin{Convertere propter servos tuos}, ``return for the sake of thy
servants''); the American psalm renders it with an adverb, so that the shared
verb is not heard in English. The refrain, by contrast, asks God to turn
\emph{us} (\latin{converte nos}).

\subsection{The second reading: a church that lacks no gift}

Paul's first letter to Corinth, to a church he had founded, opens as his letters do with a greeting and a thanksgiving. The
thanksgiving is generous: the Corinthians are rich in every word and every
knowledge, and lack no grace. The letter that follows rebukes them for
factions, for a case of incest, for lawsuits, for disorder at the Lord's
Supper and for doubting the resurrection, and calls them carnal (3:1).

The thanksgiving looks forward. The Corinthians wait ``for the manifestation
of our Lord Jesus Christ'' (1:7), the \latin{revelatio} that the \work{Ordo}'s
title puts in the mouth of the assembly. Christ ``will confirm you unto the
end without crime, in the day of the coming of our Lord'' (1:8): the Latin
\latin{sine crimine} renders a Greek word for one against whom no charge can
be brought. And ``God is faithful: by whom you are called unto the fellowship
of his Son'' (1:9).

\subsection{The Alleluia verse}

The verse is taken from Psalm 85 (84). The psalm begins by thanking God, who
``hast turned away the captivity of Jacob'' and turned from his wrath
(vv.~2--4), and then asks for mercy again: ``Shew us, O Lord,
thy mercy; and grant us thy salvation'' (v.~8). Augustine, preaching on the
psalm, began this verse with \latin{Hoc est quod cantavimus}, ``This is what
we have sung'': the verse had been sung in his church before he expounded it
(\work{Enarrationes in Psalmos} 84.9). In the Mass the verse greets the
Gospel.

\subsection{The Entrance Antiphon}

The Missal's antiphon is not a straight quotation of either Latin psalter: it has no vocative, it keeps the older
Vulgate's mocking enemies, and it says \latin{qui te exspectant}, ``who wait
for thee'', for the psalter's \latin{qui sustinent te}. The verb \latin{exspectare} recurs: the
Clementine Isaiah has God prepare good things \latin{exspectantibus te}, ``for
them that wait for thee'' (64:3; Douay 64:4); Paul's Corinthians are
\latin{exspectantibus revelationem}; and the Preface closes with those who
dare to wait (\latin{audemus exspectare}).

\subsection{The Collect}

Nearly the same prayer stands in the Old Gelasian sacramentary, a Roman
book of Mass prayers preserved in an eighth-century manuscript, as the prayer
after Communion of an Advent Mass: \latin{Da,
quaesumus, omnipotens Deus, cunctae familiae tuae hanc voluntatem in Christo
Filio tuo, Domino nostro venienti in operibus iustis aptos occurrere} (H.~A.
Wilson's edition, p.~218). The present Collect asks the same thing for ``thy
faithful'' instead of ``thy whole household'', and turns the infinitive into a
participle, \latin{occurrentes}: those who pray ask to be going out to meet
Christ with just works, so that, placed at his right hand, they may possess
the kingdom of heaven. In the judgment scene of Matthew 25 the Son of man
sets the sheep on his right hand (25:33).

The Collect asks God for a will (\latin{voluntatem}), and the works follow
from the will asked for. It names the coming one \latin{Christo tuo}, ``thy
Christ'', the Anointed who belongs to the Father, and speaks of him in the
present participle \latin{venienti}, as one already on the road.

\subsection{The Prayer over the Offerings}

This prayer is older still. It stands, word for word except for one ending,
in the Verona collection of Roman Mass prayers, among those for July (C.~L.
Feltoe's edition, p.~74, no.~372). Its argument has two members. The gifts
are offered \latin{de tuis \ldots\ collata beneficiis}, from what God has
already bestowed, so that the offering returns to God what was his. And what
God grants their devotion to accomplish in this passing time
(\latin{temporali}) is asked to become for them the reward of eternal
redemption. The prayer is said again, with the Prayer after Communion, on
Monday of the first week of Advent.

\subsection{The Preface}

The first Preface of Advent is titled \latin{De duobus adventibus Christi}, ``On
the two comings of Christ'', the two comings the \work{Universal Norms} name
in describing Advent. Its thanksgiving names the first coming
\latin{in humilitate carnis assumptae}, in the lowliness of the flesh he took,
and the second \latin{in suae gloria maiestatis}, in the glory of his
majesty. Its last clause has a history: a prayer for the Ascension in the
Verona collection asks that at the Mediator's second coming we may receive in
open gift what we now dare to \emph{hope} for (\latin{quod nunc audemus sperare
promissum}; Feltoe, p.~20). The Preface reads \latin{quod vigilantes nunc
audemus exspectare promissum}: what we now dare to \emph{await}, keeping
watch.

\subsection{The Communion Antiphon}

The antiphon is Psalm 85 (84):13 without its first word: ``the Lord will give
goodness: and our earth shall yield her fruit''. In the psalm the verse
follows the meeting of mercy and truth, justice and peace (vv.~11--12), and
truth springs out of the earth. Augustine's psalter read \latin{suavitatem},
``sweetness'', where the Vulgate has \latin{benignitatem}. The Missal uses the same antiphon
in a Mass for the sowing of the fields.

\subsection{The Prayer after Communion}

The prayer asks that the mysteries in which the assembly has shared profit
it, and gives its reason as a statement: through them God already teaches
his people, while they live among what passes, to love what is above and
cling to what lasts. Its opening and its close each stand in a different
prayer of the Verona collection, and both of those ask for what the Missal's
prayer says God already does. The opening, \latin{Prosint nobis \ldots\
frequentata mysteria}, begins a September prayer that asks that the
mysteries free those who receive them from earthly desires and teach them
(\latin{instituant}) to love heavenly things (Feltoe, p.~134). The close,
\latin{inter praetereuntia \ldots\ inhaerere mansuris}, ``among passing
things \ldots\ to cling to things that last'', ends a prayer for the
Ascension that asks God to grant his people to love heavenly things and,
among passing things, to cling already to what lasts (p.~20). Neither is an
Advent prayer.

\subsection{The solemn blessing}

The Advent blessing, which the priest may use at the end of this Mass, is
built on the same two comings as the Preface. Its three invocations ask that
those who believe the Son has come and await his coming again be sanctified;
that they be made firm in faith, glad in hope and active in charity; and that
those who rejoice in his coming in the flesh receive eternal life when he
comes again in majesty.
